\documentclass[12pt,a4paper]{article}

\usepackage{fontawesome5}
\providecommand{\faPencilAlt}{\faIcon{pencil-alt}}

\usepackage{styles/couleurs}
\usepackage{styles/mise_en_page}
\usepackage{styles/boites}
\usepackage{styles/maths}
\usepackage{styles/figures}

\usepackage{enumitem}
\usepackage{array}
\usepackage{tabularx}
\usepackage{amsmath}
\usepackage{amssymb}
\usepackage{graphicx}
\usepackage{tikz}
\usepackage{xparse}

\usetikzlibrary{calc,arrows.meta}

\setlist[itemize]{leftmargin=1.4em,itemsep=0.15em,topsep=0.2em}
\setlist[enumerate]{leftmargin=*,itemsep=0.22em,topsep=0.25em}

\newcommand{\lignepointillee}{%
  \noindent\dotfill\par
  \vspace{0.45em}
}

\newcommand{\lignesreponse}[1]{%
  \begin{tcolorbox}[
    enhanced,
    colback=white,
    colframe=NeutreC!40,
    boxrule=0.5pt,
    arc=1.5mm,
    left=3mm,
    right=3mm,
    top=2mm,
    bottom=1mm,
    before skip=2pt,
    after skip=2pt
  ]
    \foreach \i in {1,...,#1}{\lignepointillee}
  \end{tcolorbox}
}

\newcommand{\imageactivite}[3]{%
  \IfFileExists{#3}{%
    \includegraphics[width=#1,height=#2,keepaspectratio]{#3}%
  }{%
    \fbox{%
      \parbox[c][#2][c]{#1}{%
        \centering\sffamily\small
        Image à intégrer\par
        \vspace{0.3em}
        \texttt{\detokenize{#3}}
      }%
    }%
  }%
}

\newcommand{\titreseq}[1]{%
\begin{center}
  {\sffamily\LARGE\bfseries\color{TitreC}
  Séquence 5 -- Statistiques et proportionnalité}\\[0.2em]
  {\sffamily\large\bfseries\color{SavoirC}#1}
\end{center}
}

\begin{document}

\titreseq{Activité 6 -- Diagrammes circulaires et médias}

\section*{Consigne}

Observe les trois documents. Pour chacun :

\begin{enumerate}
  \item vérifie la somme des pourcentages annoncés ;
  \item compare les pourcentages et la taille apparente des secteurs ;
  \item repère les informations manquantes : titre précis, source, date,
        effectif interrogé ou unité ;
  \item décide si le diagramme peut être utilisé sans réserve.
\end{enumerate}

\begin{center}
\begin{minipage}[t]{0.31\textwidth}
\centering
\textbf{Document 1}\par\medskip
\imageactivite
  {0.95\linewidth}
  {0.23\textheight}
  {images/diagramme_media_1.png}
\end{minipage}
\hfill
\begin{minipage}[t]{0.31\textwidth}
\centering
\textbf{Document 2}\par\medskip
\imageactivite
  {0.95\linewidth}
  {0.23\textheight}
  {images/diagramme_media_2.png}
\end{minipage}
\hfill
\begin{minipage}[t]{0.31\textwidth}
\centering
\textbf{Document 3}\par\medskip
\imageactivite
  {0.95\linewidth}
  {0.23\textheight}
  {images/diagramme_media_3.png}
\end{minipage}
\end{center}

\section*{Analyse}

\renewcommand{\arraystretch}{1.45}
\begin{tabularx}{\textwidth}{
|>{\centering\arraybackslash}p{2.2cm}
|>{\centering\arraybackslash}p{2.8cm}
|>{\raggedright\arraybackslash}X
|>{\raggedright\arraybackslash}X|}
\hline
\textbf{Document} & \textbf{Somme des \%} &
\textbf{Cohérence graphique} & \textbf{Conclusion}\\
\hline
1 & & & \\
\hline
2 & & & \\
\hline
3 & & & \\
\hline
\end{tabularx}

\section*{Synthèse}

Rédige trois règles permettant de contrôler un diagramme circulaire publié dans un média.

\lignesreponse{6}

\aretenir{Un diagramme impressionnant n'est pas forcément correct.
Il faut contrôler les nombres, les secteurs et le contexte des données.}

\end{document}
